% Pursuing Stacks (cotes 134-2 to 134-8): archive cote 134-3, site page 77 (PDF page 78), typescript page 77.
% typescript. Draft compared: ps2.tex:3986-4048.
% AI-transcribed by Claude Opus 5.5 (claude-opus-5-5, adaptive thinking, effort high), Message Batches API, batch ps1 (msgbatch_01Rv5hCaNTtUt4JHx1yqoHPL), 2026-11-12; not yet checked by eye.
\scanpage{134-3/77}
\mspage{77}
only the $\hand{\Delta}_m$'s with $m\hand{\leq} n$ map into it.

Maybe I am only being emprisoned still by the preconception of fin-%
ding a homotopy interval in $\hand{\Delta^{f}}$ itself, rather than in $(\hand{\Delta^{f}})^{\wedge}$. After all,
just applying the definition of a test-category A with test-functor
i: A$\hand{\rightarrow}$ (Cat), all we have to care about is whether a) A is aspheric and
b) $i^*(\hand{\Delta}_1)$ is an aspheric element of $A^{\wedge}$. We just got to apply this to the
case of the functor
\[ i^{f} : \hand{\Delta}^{f} \hand{\longrightarrow} (\mathrm{Cat}) \]
induced by $i : \hand{\Delta} \hand{\longrightarrow}$ (Cat) above, taking into account of course the
extra trouble that $i^{f}$ is no longer fully faithful.

I just stopped \inserted[typed]{to look} with a big expectation that $\hand{\Delta}^{f}$ \uline{is} a test category
after all - but it turns out it definitely isn't ! Indeed $A_{/\hand{\Delta}_{o}\times I}$
(where $I =(i^{f})^{*}(\hand{\Delta}_1)$ ) is \inserted{again} a discrete two-point category, not aspheric. Ta-%
king the canonical functor $a \hand{\mapsto} A_{/a}$ from $A= \hand{\Delta}^{f}$ to (Cat), which is the
ultimate choice for checking whether or not A is a test category, finally
gives the answer: it is not. Because with I now the Lawvere element, we
still have that $A_{/\hand{\Delta}_{o}\times I}$ is a two-point discrete category. Thus the topos
$A^{\wedge}= (\hand{\Delta}^{f})^{\wedge}$ is'nt locally modelizing i.e. it has'nt got any homotopy \xout{xxx}\inserted[typed]{in}-%
\xout{xxxx}\inserted[typed]{terval}, which at any rate is a very big drawback I would think. The only
hope which still remains, to account for the positive features of face-%
complexes recalled above, is that $\hand{\Delta}^{f}$ is at the least a test category \uline{in}
\uline{the wide sense}\inserted{[or weak test category]}, namely that for any category C with final element, the
category $i_A j_A(C) = A_{/j_A(C)}$( = category of all \xout{objects} pairs (n,u)\inserted{,} with
u a map of the (ordered) category $A_{/\tilde{\hand{\Delta}}_n}$ \inserted{into C,}) is aspheric. The first case to
check is for C = final category $\hand{\Delta}_o$ , i.e. \xout{asyclicity} asphericity of A ,
next step would be $C =\hand{\Delta}_1$ , i.e. asphericity of the Lawvere element $L_A$
(but of course \uline{not} asphericity over the final element $e_{A^{\wedge}}$ !).

The question certainly deserves to be settled. If the answer is
affirmative i.e. $\hand{\Delta}^{f}$ \uline{is} a test category in the wide sense, then\unsure{\hand{,}} the proposi-%
\margin{-} tion stated earlier (\hand{\S}. \hand{42}), which clearly applies equally when A is a
test category in the wide sense, implies that if M is\xout{x} any full subcatego-%
ry of (Cat) containing all those C which are ``locally isomorphic'' to $\hand{\Delta}^{f}$ ,
(i.e. such that for every $x \hand{\in} C$, $C_{/x}$ is isomorphic to the ordered category
of all subsimplices of some simplex), then for the induced notion $W_M$ of
weak equivalences, M is a modelizer and the inclusion functor from M into
(Cat) is modelizing. This does not yet apply to (Ord), however it reopens
the question whether the \inserted{full sub}category of \inserted{(Ord) made up by all} \deleted{``simplicial maquettes'', which can be}
\deleted{interpreted as} ordered sets J which are locally isomorphic to $\hand{\Delta}^{f}$ in the
sense above (namely for any $x \hand{\in} J$, the \xout{xxxxx} ordered subset $J_{\hand{\leq} x}$ is isomor-%
